\begin{figure*}[t]
\centering
\setlength{\tabcolsep}{8pt}
\renewcommand{\arraystretch}{1.08}
\begin{tabularx}{\textwidth}{@{}>{\raggedright\arraybackslash}m{0.34\textwidth}>{\raggedright\arraybackslash}m{0.62\textwidth}@{}}
\fbox{%
\begin{minipage}[t][4.3cm][t]{0.31\textwidth}
\textbf{Feature-space recovery} \\
\textbf{Frozen evidence:} multiscale SAM \texttt{backbone\_fpn} features. \\
\textbf{Question:} Is texture partition structure already latent in the representation? \\
\textbf{Readout:} pooled coarsest clustering, flip-averaging, positionality diagnostics. \\
\textbf{Role in paper:} diagnostic probe only.
\end{minipage}}
&
\fbox{%
\begin{minipage}[t][4.3cm][t]{0.58\textwidth}
\textbf{Proposal-space recovery} \\
\textbf{Frozen evidence:} SAM-style mask proposals. \\
\textbf{Question:} If the bank already contains the right fragments, can a learned layer recover a coherent partition without retraining the proposal generator? \\
\textbf{Readout:} compatibility reasoning, conservative core selection, and selective repair above the frozen bank. \\
\textbf{Role in paper:} main operational contribution instantiated by \sys{}.
\end{minipage}}
\end{tabularx}
\caption{Two recoverability loci in frozen SAM-style systems. The feature route asks whether texture structure is already latent in frozen features and is used only as a diagnostic probe. The proposal route is the main operational question of the paper: how much task value can be recovered by learning commitment above a frozen proposal bank. Both routes keep the underlying segmenter frozen.}
\label{fig:recoverability_loci}
\end{figure*}
